\chapter{Технологический раздел}
\section{Средства реализации}
Алгоритмы для данной лабораторной работы были реализованы на языке C++ при использовании компилятора gcc версии 13.3.1, так как в стандартной библиотеке приведённого языка присутствуют все необходимые структуры данных для реализации требуемых алгоритмов решения задачи коммивояжёра.

\section{Реализация алгоритмов}
Реализация муравьиного алгоритма и алгоритма полного перебора для решения задачи коммивояжёра приведена в приложении А.

\section{Тестирование}

Модульные тесты рассмотрены в таблице~\ref{tbl:test}. Все тесты были успешно пройдены.

\begin{center}
	\captionsetup{justification=raggedright,singlelinecheck=off}
	\begin{longtable}[c]{|c|c|c|c|c|}
		\caption{Модульные тесты\label{tbl:test}} \\ \hline
		Матрица смежности & Ожидаемый результат & Результат программы \\
		\hline
		$ \begin{pmatrix}
			0 & 1 & 2 & 3 & 4 \\
			1 & 0 & 5 & 6 & 7 \\
			2 & 5 & 0 & 8 & 9 \\
			3 & 6 & 8 & 0 & 10 \\
			4 & 7 & 9 & 10 & 0 
		\end{pmatrix}$ &
		27, [0, 2, 1, 3, 4] &
		27, [0, 2, 1, 3, 4] \\
		\hline
		$ \begin{pmatrix}
			0 & 1 & 2 & 3 \\
			1 & 0 & 4 & 5 \\
			2 & 4 & 0 & 6 \\
			3 & 5 & 6 & 0 \\
		\end{pmatrix}$ &
		14, [0, 1, 2, 3] &
		14, [0, 1, 2, 3] \\
		\hline
		$ \begin{pmatrix}
			0 & 27 & 12 & 21 \\
			27 & 0 & 16 & 24 \\
			12 & 16 & 0 & 11 \\
			21 & 24 & 11 & 0 
		\end{pmatrix}$ &
		73, [0, 2, 1, 3] &
		73, [0, 2, 1, 3] \\
		\hline
	\end{longtable}
\end{center}